\section{模型评价与推广}

\subsection{模型优点}
四问从单次测向推进到多源闭环搜索：先建立位置约束，再讨论下一检测点，最后同时处理已发现源的清除与未知频道的发现。各阶段可以核验的内容如下：
\begin{enumerate}
  \item 问题一把示向度误差写成半平面交，定位区域的每条边都能追溯到具体观测；直径与圆覆盖分开检验，避免把二者混同；
  \item 问题二以可行域投影和双侧偏移给出明确的候选区域，前移比例与侧移距离分别控制纵向行程和交会几何，并与后两问的首次复测规则一致；
  \item 问题三按频道维护保守可行域，并以已完成扫描点的覆盖证书判断未知频道是否得到全域排查；
  \item 问题四不以单点无信号排除定向源，而通过多方位观测建立对任意发射朝向成立的阴性证书。
\end{enumerate}

\subsection{模型局限}
当前模型与证据仍有以下限制：
\begin{enumerate}
  \item 问题二的算例只检验一个事后设定的源位置；双侧候选构造不是全域最优性证明，实际接收与定位效果仍须依据后续观测及三、四问测试评价；
  \item 发现覆盖和有限网格兜底分别给出条件性几何保证；若动作数或真实运行时间先耗尽，不能据此声称已清除全部干扰源；
  \item 当前三、四问缺少与选定程序对应的演练和正式测试数值，路线优化与多方位补测的实际时间收益仍需日志验证；
  \item 通信和转向耗时未计入当前时间模型，若用于实体设备，应先测量这些开销并重新标定决策参数。
\end{enumerate}

\subsection{模型推广}
若干扰源可能移动，静态扇形交集需随新观测实时更新；若场地存在障碍，第二检测点的直线移动距离应改为可通行路径长度。多机器狗协同时，还需分配检测区域与目标，避免重复测向和路径冲突。扩展时仍可沿用“记录观测、更新可行域、比较下一行动”的闭环；能否保证接收，则须按全向或定向场景分别判断。
